\section{Conclusion} \label{sec:Further_Steps}

This review has examined the transit photometry detection method and the challenges associated with its use. Various issues such as instrumental noise and stellar variability impact the ability for astronomers to detect potential exoplanets. To mitigate these difficulties, several space missions have been developed to monitor as many possible stars, using preprocessing steps to improve detectability.

\subsection{Further Steps}

Moving forward, machine learning is likely to play an increasingly important role in the field. Modern surveys have collected hundreds of thousands of light curves, making the analysis of all available data very impractical. With increasing dataset sizes and the continued challenge of false positives, convolutional neural networks (CNNs) provide a valuable addition to the research process. The challenges associated with transit detection and identification of astrophysical false positives remains central to the developing automated detection pipelines. As photometric surveys continue to expand in scale, CNNs can analyse large volumes of data with great efficiency and identify potential exoplanet candidates. Continued development of these techniques will be essential for improving the scientific return from current and future missions such as TESS and PLATO.

The \textit{Kepler} mission provides an extensive archive of high-precision stellar light curves, making it an ideal starting point for developing and testing machine learning approaches. Its large catalogue of classified signals offers a strong foundation for training and evaluating models before extending these techniques to TESS and PLATO.

\subsection{Proposed Project}

The question is whether a CNN can be trained to correctly classify light curves with minimal preprocessing in order to reduce computational cost and time. Current research often trains CNN models on cleaned and phase-folded light curves, where preprocessing makes the relevant transit features easier to identify. When managing hundreds of thousands of light curves, this preprocessing becomes highly time consuming. CNN pipelines are relatively straightforward to set up, reducing preprocessing time and allowing more time spent on analysis. However, this preprocessing is conducted to reduce the false positive rate leading to improved classification accuracy. The overall goal is to distinguish true planetary transits from common false positives and understand if speed can be gained without loss in quality.  

To apply this strategy of handling raw light curves, the first step is to establish a working CNN pipeline that runs without error. This begins with a simple case, making sure accuracy and other metrics can be calculated and plotted over time. To do this, a simple process such as uniform randomness is applied to introduce white noise to the data. The goal of this dataset is to represent cases of transits, eclipsing binaries, and normal stars. Using this synthetic data, the algorithm is trained to learn basic features of different light curve classes showing that the neural network is determining meaningful patterns through its feature maps. Once this model is developed, it works as a baseline for building into more complex algorithms. 

Real data is not nearly as clean and so this simple trained model will not be sufficient enough for handling realistic data. Once the first model works, the next step is the implementation of more realistic synthetic data. One effective way to do this is to overlay real stellar noise over the earlier synthetic data. This immediately makes the classification task more similar to real light curves, as the embedded features are now surrounded by realistic noise. This allows the CNN parameters to be adjusted to make is easier for the model to work with the data. 

This harder model must now look through the noise and detect the underlying pattern. Once this model is prepared, the next step is to move on to real data. \textit{Kepler} is the target domain for evaluating whether the earlier learning transfers. The \textit{Kepler} catalogue contains a large selection of confirmed and potential exoplanets, as well as eclipsing binaries to use in modelling. Using these, it is straightforward to build a model focused on the imperfections present in real light curves. Further adjustments can then be made to the training process to better analyse model performance and determine the most effective parameters for light curve classifications.

Overall, the training process is divided across three stages, creating a progressive approach. Each stage of learning is intended to prepare the CNN for the next. The final model is therefore expected to combine the basic feature learning of the first stage, the robustness of realistic variability and adaptation to real \textit{Kepler} data. This method allows the model to focus on one process at a time. Success is determined by whether this staged approach maintains strong classification performance throughout the training process. This would show that performance quality can be maintained while reducing preprocessing. 
